
\subsection{算子 \texorpdfstring{$\mathfrak q$}{q} 的构造纲要}
\label{subsec:3-8-6}
\liftIndex{I356}\liftIndexRangeStart{I358}

定理 3.8.32 的证明有一个微妙之处，即光滑奇异单形可数集
$\mathcal X_{1+}(L)$ 的选择。我们把这项选择推迟到第 7.4 节；那里
以第 7.2 节的讨论为基础。本小节先假设存在这样的
$\mathcal X_{1+}(L)$，概述算子 $\mathfrak q$ 的构造以及
\eqref{eq:3-8-33}、\eqref{eq:3-8-34-1} 和
\eqref{eq:3-8-34-2} 的证明。

考虑第 2.1.2 小节引入的模空间
$\mathcal M_{k+1,\ell}^{\mathrm{main}}(L,\beta)$。
把 $k+1$ 个边界标记点记作 $z_0,\ldots,z_k$，把 $\ell$ 个内部标记点
记作 $z_1^+,\ldots,z_\ell^+$，并置
\[
 \vec z=(z_1,\ldots,z_k),
 \qquad
 \vec z^+=(z_1^+,\ldots,z_\ell^+).
\]
评价映射
\[
\begin{aligned}
 ev=(ev_{\mathrm{int}},ev_\partial,ev_0)\colon
 \mathcal M_{k+1,\ell}^{\mathrm{main}}(L,\beta)
 \longrightarrow M^\ell\times L^{k+1}
\end{aligned}
\]
定义为
\[
\begin{aligned}
 &ev\bigl(w,(\vec z^+,\vec z,z_0)\bigr)\\
 &\quad=
 \bigl(
 w(z_1^+),\ldots,w(z_\ell^+),
 w(z_1),\ldots,w(z_k),w(z_0)
 \bigr).
\end{aligned}
\]
给定
\[
 Q_i\in\mathcal X(M)\quad(i=1,\ldots,\ell),
 \qquad
 P_i\in\mathcal X_{1+}(L)\quad(i=1,\ldots,k),
\]
定义纤维积
\[
\begin{aligned}
 &\mathcal M_{k+1,\ell}^{\mathrm{main}}
 (L,\beta;\vec Q,\vec P)\\
 &\quad=
 \mathcal M_{k+1,\ell}^{\mathrm{main}}(L,\beta)
 _{(ev_{\mathrm{int}},ev_\partial)}
 \mathbin{\times_{M^\ell\times L^k}}
 (Q_1\times\cdots\times Q_\ell)
 \times(P_1\times\cdots\times P_k),
\end{aligned}
\]
其中
\[
 \vec Q=(Q_1,\ldots,Q_\ell),
 \qquad
 \vec P=(P_1,\ldots,P_k).
\]
下面两个命题将分别在第 7.1.4 小节和第 7.4 节证明。

\begin{proposition}{3.8.65}\label{prop:3-8-65}
$\mathcal M_{k+1,\ell}^{\mathrm{main}}
(L,\beta;\vec Q,\vec P)$ 带有仓西结构；在边界上，这一结构与
\[
 \mathcal M_{k'+1}^{\mathrm{main}}(L,\beta';\vec P')
\]
（见 \eqref{eq:3-5-1}）以及
\[
 \mathcal M_{k'+1,\ell'}^{\mathrm{main}}
 (L,\beta';\vec Q',\vec P')
\]
的仓西结构相容。
\end{proposition}

\begin{proposition}{3.8.66}\label{prop:3-8-66}
可以选择 $\mathcal X_{1+}(L)$，使得存在一族多值截面
$\mathfrak s_{\beta,\vec Q,\vec P}$：它与
$\mathfrak s_{\beta',\vec P'}$ 相容，在端部彼此相容，并且满足
\[
 \bigl(
 \mathfrak s_{\beta,\vec Q,\vec P}^{-1}(0),ev_0
 \bigr)
 \in C_{1+}(L).
\]
\end{proposition}

这里取适当的三角剖分并给它的顶点排序，从而把
\[
 \bigl(
 \mathfrak s_{\beta,\vec Q,\vec P}^{-1}(0),ev_0
 \bigr)
\]
视为光滑奇异链。还要指出，在命题 3.8.65 的证明中，对
$P_i\in\mathcal X_{1+}(L)\setminus\mathcal X_1(L)$ 也必须使用这样的
$\mathfrak s_{\beta,\vec P}$。记
\[
 \mathcal M_{k+1,\ell}^{\mathrm{main}}
 (L,\beta;\vec Q,\vec P)^{\mathfrak s}
 =
 \bigl(
 \mathfrak s_{\beta,\vec Q,\vec P}^{-1}(0),ev_0
 \bigr).
\]

\begin{lemma}{3.8.67}\label{lem:3-8-67}
设 $R\supseteq\mathbb Q$。可以选择
$\mathfrak s_{\beta,\vec Q,\vec P}$，使对每个
$\sigma\in\mathfrak S_\ell$ 都有
\[
\begin{aligned}
 &\mathcal M_{k+1,\ell}^{\mathrm{main}}
 (L,\beta;\sigma(\vec Q),\vec P)
 ^{\mathfrak s_{\beta,\sigma(\vec Q),\vec P}}\\
 &\qquad=
 (-1)^{\epsilon(\sigma)}
 \mathcal M_{k+1,\ell}^{\mathrm{main}}
 (L,\beta;\vec Q,\vec P)
 ^{\mathfrak s_{\beta,\vec Q,\vec P}},
\end{aligned}
\]
其中
\[
\begin{aligned}
 \sigma(\vec Q)
 &=
 (Q_{\sigma(1)},\ldots,Q_{\sigma(\ell)}),\\
 \epsilon(\sigma)
 &=
 \sum_{\substack{i<j\\\sigma(i)>\sigma(j)}}
 \deg Q_i\,\deg Q_j.
\end{aligned}
\]
\end{lemma}

现在定义
\begin{equation}
\left\{
\begin{aligned}
 \mathfrak q_{\ell,k,\beta}
 \bigl(
 (Q_1,\ldots,Q_\ell)\otimes(P_1,\ldots,P_k)
 \bigr)
 &:=
 \frac1{\ell!}
 \mathcal M_{k+1,\ell}^{\mathrm{main}}
 (L,\beta;\vec Q,\vec P)^{\mathfrak s},\\
 \mathfrak q_{\ell,k}
 &:=
 \sum_{\beta\in\pi_2(M,L)}
 T^{\omega[\beta]}e^{\mu_L(\beta)/2}
 \mathfrak q_{\ell,k,\beta}.
\end{aligned}
\right.
\tag{3.8.68}\label{eq:3-8-68}
\end{equation}
\liftIndexRangeEnd{I358}

下面验证 \eqref{eq:3-8-33}。为此研究
$\mathcal M_{k+1,\ell}^{\mathrm{main}}
(L,\beta;\vec Q,\vec P)$ 的边界。设
\[
 \left(
 \bigl(w^{(m)},\vec z^{+(m)},\vec z^{(m)},z_0^{(m)}\bigr),
 \bigl(\vec p_{\mathrm{int}}^{(m)},\vec p_\partial^{(m)}\bigr)
 \right)
\]
是其中的一个发散序列，其中
\[
\begin{aligned}
 \vec p_{\mathrm{int}}^{(m)}
 &=
 (p_{\mathrm{int},1}^{(m)},\ldots,
  p_{\mathrm{int},\ell}^{(m)}),
 &p_{\mathrm{int},i}^{(m)}&\in Q_i,\\
 \vec p_\partial^{(m)}
 &=
 (p_{\partial,1}^{(m)},\ldots,p_{\partial,k}^{(m)}),
 &p_{\partial,i}^{(m)}&\in P_i.
\end{aligned}
\]
有以下三种情形：

\noindent\textup{(3.8.69.1)}
$\displaystyle
\lim_{m\to\infty}p_{\partial,i}^{(m)}\in\partial P_i$。

\noindent\textup{(3.8.69.2)}
$\displaystyle
\lim_{m\to\infty}p_{\mathrm{int},i}^{(m)}\in\partial Q_i$。

\noindent\textup{(3.8.69.3)}
$w^{(m)}\colon\Sigma^{(m)}\to M$ 收敛到映射
\[
 w^{(\infty)}\colon\Sigma_1\cup\Sigma_2\longrightarrow M,
\]
其中 $\Sigma_1,\Sigma_2$ 是沿边界粘合的两个亏格零有边黎曼曲面。
见图 3.8.4。

\begin{figure}[htbp]
  \centering
  \setlength{\unitlength}{0.92mm}
  \begin{picture}(126,58)(0,-4)
    \put(42,29){\circle{36}}
    \put(78,29){\circle{36}}
    \put(27,39){\line(1,-1){2.4}}
    \put(24,29){\line(1,0){2.4}}
    \put(29,16){\line(1,1){2.4}}
    \put(48,13){\line(-1,1){2.4}}
    \put(23,43){$P$}
    \put(17,27){$P$}
    \put(24,9){$P$}
    \put(48,6){$P$}
    \put(67,43){\line(1,-1){2.4}}
    \put(77,47){\line(0,-1){2.4}}
    \put(91,40){\line(-1,-1){2.4}}
    \put(88,14){\line(-1,1){2.4}}
    \put(65,47){$P$}
    \put(76,51){$P$}
    \put(92,43){$P$}
    \put(88,7){$P$}
    \put(96,29){\vector(1,0){9}}
    \put(107,27){$ev_0$}
    \put(37,28){\circle*{1.5}}
    \put(49,29){\circle*{1.5}}
    \put(38,29){\line(0,1){7}}
    \put(50,30){\line(0,1){7}}
    \put(36,38){$Q$}
    \put(49,39){$Q$}
    \put(70,27){\circle*{1.5}}
    \put(81,35){\circle*{1.5}}
    \put(84,21){\circle*{1.5}}
    \put(70,28){\line(0,1){7}}
    \put(81,36){\line(0,1){7}}
    \put(84,22){\line(0,1){7}}
    \put(68,37){$Q$}
    \put(80,45){$Q$}
    \put(84,31){$Q$}
  \end{picture}
  \par\smallskip
  \renewcommand{\thefigure}{3.8.4}
  \refstepcounter{figure}
  \textbf{\figurename~\thefigure}\label{fig:3-8-4}
\end{figure}

考虑情形 (3.8.69.3)\translatornote{原书在引理 3.8.67 中将
$\sigma(\vec Q)$ 的末项误写成 $Q_{\sigma(k)}$；因
$\sigma\in\mathfrak S_\ell$ 且 $\vec Q$ 长为 $\ell$，译文恢复为
$Q_{\sigma(\ell)}$。}。可取 $\Sigma_1,\Sigma_2$ 使
$z_0^{(\infty)}\in\Sigma_2$。对每个划分
\[
 A_1\cup A_2=\{1,\ldots,\ell\},
 \qquad A_1\cap A_2=\varnothing,
\]
都有 (3.8.69.3) 的相应分支，使
$z_i^{+(\infty)}\in\Sigma_j$ 当且仅当 $i\in A_j$。
\translatornote{原书这里写成“$Q_i\in\Sigma_j$”，但 $Q_i$ 是
$M$ 中的链，不是定义域曲面上的点；划分所记录的实际是第 $i$ 个
内部标记点落在哪个分支。译文依紧前退化序列恢复
$z_i^{+(\infty)}\in\Sigma_j$。}
把这些分支记作
\[
 \partial_{A_1,A_2}
 \mathcal M_{k+1,\ell}^{\mathrm{main}}
 (L,\beta;\vec Q,\vec P).
\]

对 $i\leq j$，置
\[
\begin{aligned}
 \vec P_{i,j}&=(P_i,\ldots,P_j),\\
 \vec P_{\leq i}&=(P_1,\ldots,P_i),\\
 \vec P_{>j}&=(P_{j+1},\ldots,P_k).
\end{aligned}
\]
将 $A_1$ 的元素依次写作
$m_{A_1,1}<\cdots<m_{A_1,\#A_1}$，并以相同方式定义
$m_{A_2,i}$。再置
\[
\begin{aligned}
 \sigma(\vec Q;A_j)_i
 &:=
 Q_{\sigma(m_{A_j,i})},\\
 \sigma(\vec Q;A_j)
 &:=
 \bigl(
 Q_{\sigma(m_{A_j,1})},\ldots,
 Q_{\sigma(m_{A_j,\#A_j})}
 \bigr),\\
 \vec Q_{A_j}
 &:=
 (Q_{m_{A_j,1}},\ldots,Q_{m_{A_j,\#A_j}}).
\end{aligned}
\]
于是\translatornote{原书在本式以及下面的 (3.8.70) 中都把曲线类条件
写成 $\beta_j\ne0$，但求和只定义了 $\beta_1,\beta_2$，而 $j$ 已是
边界输入位置。内层气泡模空间的曲线类为 $\beta_1$；译文据这一型别
恢复 $\beta_1\ne0$。}
\[
\begin{aligned}
&\frac1{\ell!}
 \sum_{(A_1,A_2)}
 \sum_{\sigma\in\mathfrak S_\ell}
 (-1)^{\epsilon(\sigma)}
 \partial_{A_1,A_2}
 \mathcal M_{k+1,\ell}^{\mathrm{main}}
 (L,\beta;\sigma(\vec Q),\vec P)^{\mathfrak s}\\
&\quad=
 \frac1{\ell!}
 \sum_{\sigma\in\mathfrak S_\ell}
 \sum_{\substack{\beta_1+\beta_2=\beta\\\beta_1\ne0}}
 \sum_{(A_1,A_2)}
 \sum_{0\leq i\leq j\leq k}
 (-1)^{\epsilon(\sigma)}\\
&\qquad\quad\cdot
 \mathcal M_{k-j+i,\#A_2}^{\mathrm{main}}
 \Bigl(
 L;\beta_2;\sigma(\vec Q;A_2);
 \vec P_{\leq i},\\
&\hspace{53mm}
 \mathcal M_{j-i+1,\#A_1}^{\mathrm{main}}
 (L,\beta_1;\sigma(\vec Q;A_1);\vec P_{i,j})^{\mathfrak s},
 \vec P_{>j}
 \Bigr).
\end{aligned}
\]
利用引理 3.8.67，可见上式等于
\[
\begin{aligned}
&\sum_{(A_1,A_2)}
 \sum_{\substack{\beta_1+\beta_2=\beta\\\beta_1\ne0}}
 \sum_{0\leq i\leq j\leq k}
 \sum_{\sigma_1\in\mathfrak S_{\#A_1}}
 \sum_{\sigma_2\in\mathfrak S_{\#A_2}}
 \frac{(-1)^{\epsilon_{i,j,A_1,A_2}}}
 {\#A_1!\,\#A_2!}\\
&\quad\cdot
 \mathcal M_{k-j+i,\#A_2}^{\mathrm{main}}
 \Bigl(
 L;\beta_2;\sigma_2(\vec Q_{A_2});
 \vec P_{\leq i},\\
&\hspace{48mm}
 \mathcal M_{j-i+1,\#A_1}^{\mathrm{main}}
 (L,\beta_1;\sigma_1(\vec Q_{A_1});\vec P_{i,j})^{\mathfrak s},
 \vec P_{>j}
 \Bigr).
\end{aligned}
\]
对 $\vec Q=(Q_1,\ldots,Q_\ell)$，置
\[
 [\vec Q]
 =
 \sum_{\sigma\in\mathfrak S_\ell}\sigma(\vec Q).
\]
于是\translatornote{原书在 (3.8.70) 的外层算子下标中，于
$\#A_2$ 与 $k-j+i$ 之间多出一个未定义的 $\ell_2$；而
\eqref{eq:3-8-68} 的分曲线类算子只有 $(\ell,k,\beta)$ 三个下标，
内层算子也遵守这一型别。译文删除该游离下标。}
\begin{equation}
\begin{aligned}
&\frac1{\ell!}
 \sum_{(A_1,A_2)}
 \sum_{\sigma\in\mathfrak S_\ell}
 (-1)^{\epsilon(\sigma)}
 \partial_{A_1,A_2}
 \mathcal M_{k+1,\ell}^{\mathrm{main}}
 (L,\beta;\sigma(\vec Q),\vec P)^{\mathfrak s}\\
&=
 \sum_{(A_1,A_2)}
 \sum_{\substack{\beta_1+\beta_2=\beta\\\beta_1\ne0}}
 \sum_{0\leq i\leq j\leq k}
 \mathfrak q_{\#A_2,k-j+i,\beta_2}
 \Bigl(
 [\vec Q_{A_2}]\otimes
 \vec P_{\leq i},\\
&\hspace{42mm}
 \mathfrak q_{\#A_1,j-i+1,\beta_1}
 \bigl([\vec Q_{A_1}]\otimes\vec P_{i,j}\bigr),
 \vec P_{>j}
 \Bigr).
\end{aligned}
\tag{3.8.70}\label{eq:3-8-70}
\end{equation}
注意
\[
 \Delta\vec Q
 =
 \sum_{A_1<A_2}\vec Q_{A_1}\otimes\vec Q_{A_2},
\]
其中 $A_1<A_2$ 表示对任意 $i\in A_1,j\in A_2$ 都有 $i<j$。
因此
\[
 \Delta[\vec Q]
 =
 \sum_{(A_1,A_2)}
 [\vec Q_{A_1}]\otimes[\vec Q_{A_2}].
\]
所以 \eqref{eq:3-8-70} 给出 \eqref{eq:3-8-33} 中不含经典边界算子
的各项。涉及经典边界算子的各项，或者来自 (3.8.69.1)、
(3.8.69.2)，或者来自
\[
 \partial\mathcal M_{k+1,\ell}^{\mathrm{main}}
 (L,\beta;\vec Q,\vec P).
\]
至此 \eqref{eq:3-8-33} 得证。

下面证明 (3.8.34)。式 \eqref{eq:3-8-34-2} 由构造立即可见。为证明
\eqref{eq:3-8-34-1}，注意
$\mathcal M_{k+1,1}^{\mathrm{main}}(L,\beta_0)$ 由取值于 $L$ 的常值
映射组成；回忆 $\beta_0=0\in\pi_2(M,L)$。因此，当 $k=0$ 时，
\[
 \mathcal M_{1,1}^{\mathrm{main}}(L,\beta_0,Q)
 =
 L\cap Q.
\]
这就得到 \eqref{eq:3-8-34-1} 中 $k=0,\ell=1$ 的情形。当
$k>0,\ell=1$ 时，
$\mathcal M_{k+1,1}^{\mathrm{main}}(L,\beta_0;Q,\vec P)$ 的支集包含于
\[
 \sum_i
 \mathfrak m_{k+1}^L
 (P_1,\ldots,P_i,L\cap Q,P_{i+1},\ldots,P_k).
\]
另一方面，
\[
\begin{aligned}
 &\deg
 \mathfrak m_{k+1}^L
 (P_1,\ldots,P_i,L\cap Q,P_{i+1},\ldots,P_k)\\
 &\qquad=
 \deg\mathfrak q_{1,k}
 \bigl(Q\otimes(P_1,\ldots,P_k)\bigr)+1.
\end{aligned}
\]
故置
\[
 \mathfrak q_{1,k,\beta_0}
 \bigl(Q\otimes(P_1,\ldots,P_k)\bigr)=0.
\]
$\ell>1$ 的情形也可同样处理；见注 7.2.172 (3)。

(3.8.35) 与 (3.8.36) 的证明，是把上述论证与第 7.3 节的论证合并；
见第 7.3 节末。

\subsection{算子 \texorpdfstring{$\mathfrak r$}{r} 与
\texorpdfstring{$A_\infty$}{A-infinity} 双模}
\label{subsec:3-8-7}
\liftIndex{I199}\liftIndex{I200}

为了证明定理 I，需要把第 3.8.5 小节的构造同滤过 $A_\infty$ 双模
结合起来。设 $(L^{(0)},L^{(1)})$ 是 $M$ 中的一对相对 spin
拉格朗日子流形。像第 3.8.5 小节那样，选取
\[
 \mathcal X_{1+}(L^{(1)}),\qquad
 \mathcal X_{1+}(L^{(0)}),\qquad
 \mathcal X(M).
\]
这里对 $L^{(0)}$ 和 $L^{(1)}$ 使用同一个 $\mathcal X(M)$。假设
$L^{(0)}$ 与 $L^{(1)}$ 干净相交，并分解
\[
 L^{(1)}\cap L^{(0)}=\bigcup_h R_h.
\]
只须考虑博特--莫尔斯情形，因为它包含横截情形。对每个 $R_h$ 取链的
可数集 $\mathcal X_1(R_h)$，像第 3.7.5 小节那样定义
$C(R_h;\Lambda_{0,\mathrm{nov}})$，并仍按该小节置
$C^k(L^{(1)},L^{(0)};\Lambda_{0,\mathrm{nov}})$。现在有：

\begin{theorem}{3.8.71}\label{thm:3-8-71}
存在光滑奇异单形的可数集
\[
 \mathcal X_{1+}(R_h)\supset\mathcal X_1(R_h),
\]
它生成自由 $\Lambda_{0,\mathrm{nov}}$-模
$C_{1+}(L^{(1)},L^{(0)};\Lambda_{0,\mathrm{nov}})$；第 3.7.5 小节
构造的滤过 $A_\infty$ 双模结构可以扩张到这个自由模上。并且存在
一族次数为 $+1$ 的算子
\[
\begin{aligned}
 \mathfrak r_{\ell;k_1,k_0}\colon&
 E_\ell C(M;\Lambda_{0,\mathrm{nov}})[2]
 \otimes
 B_{k_1}C_{1+}(L^{(1)};\Lambda_{0,\mathrm{nov}})[1]\\
 &\otimes
 C_{1+}(L^{(1)},L^{(0)};\Lambda_{0,\mathrm{nov}})[1]\\
 &\otimes
 B_{k_0}C_{1+}(L^{(0)};\Lambda_{0,\mathrm{nov}})[1]\\
 &\longrightarrow
 C_{1+}(L^{(1)},L^{(0)};\Lambda_{0,\mathrm{nov}})[1],
\end{aligned}
\]
满足
\begin{equation}
\begin{aligned}
0={}&
 \mathfrak r_{\ell;k_1,k_0}
 \bigl(
 \widehat\delta_M(\mathbf y)
 \otimes\mathbf x\otimes v\otimes\mathbf z
 \bigr)\\
&+
 \sum_{c_1,c_2}
 (-1)^{
  \deg y_{c_1}^{(2;2)}
   \deg'\mathbf x_{c_2}^{(3;1)}
  +\deg'\mathbf x_{c_2}^{(3;1)}
  +\deg y_{c_1}^{(2;1)}}\\
&\qquad\cdot
 \mathfrak r\Bigl(
 y_{c_1}^{(2;1)}
 \otimes
 \bigl(
 \mathbf x_{c_2}^{(3;1)}
 \otimes
 \mathfrak q(
 y_{c_1}^{(2;2)}
 \otimes\mathbf x_{c_2}^{(3;2)})
 \otimes\mathbf x_{c_2}^{(3;3)}
 \bigr)
 \otimes v\otimes\mathbf z
 \Bigr)\\
&+
 \sum_{c_1,c_2,c_3}
 (-1)^{
  \deg y_{c_1}^{(2;2)}
   \deg'\mathbf x_{c_2}^{(2;1)}
  +\deg'\mathbf x_{c_2}^{(2;1)}
  +\deg y_{c_1}^{(2;1)}}\\
&\qquad\cdot
 \mathfrak r\Bigl(
 y_{c_1}^{(2;1)}
 \otimes\mathbf x_{c_2}^{(2;1)}
 \otimes
 \mathfrak r\bigl(
 y_{c_1}^{(2;2)}
 \otimes\mathbf x_{c_2}^{(2;2)}
 \otimes v\otimes\mathbf z_{c_3}^{(2;1)}
 \bigr)
 \otimes\mathbf z_{c_3}^{(2;2)}
 \Bigr)\\
&+
 \sum_{c_1,c_3}
 (-1)^{
  (\deg y_{c_1}^{(2;2)}+1)
  (\deg'\mathbf x+\deg'v+
   \deg'\mathbf z_{c_3}^{(3;1)})
  +\deg y_{c_1}^{(2;1)}}\\
&\qquad\cdot
 \mathfrak r\Bigl(
 y_{c_1}^{(2;1)}
 \otimes
 \bigl(
 \mathbf x\otimes v\otimes
 \bigl(
 \mathbf z_{c_3}^{(3;1)}
 \otimes
 \mathfrak q(
 y_{c_1}^{(2;2)}
 \otimes\mathbf z_{c_3}^{(3;2)})
 \otimes\mathbf z_{c_3}^{(3;3)}
 \bigr)
 \bigr)
 \Bigr).
\end{aligned}
\tag{3.8.72}\label{eq:3-8-72}
\end{equation}
这里 $\widehat\delta_M$ 定义于第 3.8.4 小节。此外还有

\noindent\textup{(3.8.73.1)}
\[
 \mathfrak r_{0;k_1,k_0}=\mathfrak n_{k_1,k_0}.
\]

\noindent\textup{(3.8.73.2)}
\[
 \mathfrak r_{\ell;k_1,k_0}
 \equiv0\pmod{\Lambda^+_{0,\mathrm{nov}}}
 \qquad(\ell\ne0).
\]

还可以把 $\mathfrak r$ 扩张为
\[
\begin{aligned}
 \mathfrak r^+_{\ell;k_1,k_0}\colon&
 E_\ell C(M;\Lambda_{0,\mathrm{nov}})[2]
 \otimes
 B_{k_1}C_{1+}(L^{(1)};\Lambda_{0,\mathrm{nov}})^+[1]\\
 &\otimes
 C_{1+}(L^{(1)},L^{(0)};\Lambda_{0,\mathrm{nov}})[1]\\
 &\otimes
 B_{k_0}C_{1+}(L^{(0)};\Lambda_{0,\mathrm{nov}})^+[1]\\
 &\longrightarrow
 C_{1+}(L^{(1)},L^{(0)};\Lambda_{0,\mathrm{nov}})[1].
\end{aligned}
\]
把 \eqref{eq:3-8-72} 中的 $\mathfrak r,\mathfrak q$ 分别换成
$\mathfrak r^+,\mathfrak q^+$，它仍满足同一公式。它还满足
\[
\begin{aligned}
 &\mathfrak r^+_{\ell;k_1,k_0}
 (\,\mathbf y\otimes x_1\otimes\cdots\otimes
 \mathbf e_1^+\otimes\cdots\otimes x_{k_1-1}
 \otimes v\otimes z_1\otimes\cdots\otimes z_{k_0}\,)
 =0,\\
 &\mathfrak r^+_{\ell;k_1,k_0}
 (\,\mathbf y\otimes x_1\otimes\cdots\otimes x_{k_1}
 \otimes v\otimes z_1\otimes\cdots\otimes
 \mathbf e_0^+\otimes\cdots\otimes z_{k_0-1}\,)
 =0
\end{aligned}
\]
（$k_1+k_0\geq2$）。这里 $\mathbf e_i^+$ 是
\[
\begin{aligned}
 C_{1+}(L^{(i)};\Lambda_{0,\mathrm{nov}})^+
 ={}&
 C_{1+}(L^{(i)};\Lambda_{0,\mathrm{nov}})
 \oplus\Lambda_{0,\mathrm{nov}}\mathbf e_i^+\\
 &\oplus\Lambda_{0,\mathrm{nov}}\mathbf f_i
\end{aligned}
\]
中的幺元，与第 3.3 节相同。此外，
\[
 \mathfrak r^+_{0;k_1,k_0}
 =
 \mathfrak n^+_{k_1,k_0},
\]
其中 $\mathfrak n^+_{k_1,k_0}$ 是
$C_+(L^{(1)},L^{(0)};\Lambda_{0,\mathrm{nov}})$ 上的滤过双模结构；
其左、右代数分别为
\[
 C_{1+}(L^{(1)};\Lambda_{0,\mathrm{nov}})^+,
 \qquad
 C_{1+}(L^{(0)};\Lambda_{0,\mathrm{nov}})^+.
\]
\end{theorem}

下一小节将给出证明纲要；完整证明见第 7.4.8--7.4.9 小节。在那里，
我们用 $L_\infty$ 模同态的语言重新表述定理 3.8.71，特别是
\eqref{eq:3-8-72}；见定理 7.4.154。

现在取
\[
 \mathfrak b\in C^2(M;\Lambda^+_{0,\mathrm{nov}}),
 \qquad
 \delta_M\mathfrak b=0.
\]
定义
\[
\begin{aligned}
 \mathfrak n^{\mathfrak b}_{k_1,k_0}\colon&
 B_{k_1}C_{1+}(L^{(1)};\Lambda_{0,\mathrm{nov}})[1]
 \otimes C_{1+}(L^{(1)},L^{(0)};\Lambda_{0,\mathrm{nov}})\\
 &\otimes
 B_{k_0}C_{1+}(L^{(0)};\Lambda_{0,\mathrm{nov}})[1]\\
 &\longrightarrow
 C_{1+}(L^{(1)},L^{(0)};\Lambda_{0,\mathrm{nov}})
\end{aligned}
\]
如下：
\[
 \mathfrak n^{\mathfrak b}_{k_1,k_0}
 (\mathbf x\otimes v\otimes\mathbf z)
 =
 \mathfrak r(
 e^{\mathfrak b}\otimes\mathbf x\otimes v\otimes\mathbf z
 ).
\]

进一步，像定义--引理 3.7.13 那样，对满足
\[
 b_i\in C_{1+}(L^{(i)};\Lambda_{0,\mathrm{nov}}),
 \qquad
 b_i\equiv0\pmod{\Lambda^+_{0,\mathrm{nov}}},
\]
的 $b_i$，定义
\[
\begin{aligned}
 &{}^{(\mathfrak b,b_1)}
 \mathfrak n^{(\mathfrak b,b_0)}
 (\mathbf x\otimes v\otimes\mathbf z)\\
 &\quad=
 \mathfrak r\bigl(
 e^{\mathfrak b}\otimes
 e^{b_1}\otimes x_1\otimes e^{b_1}\cdots
 e^{b_1}\otimes x_{k_1}\otimes e^{b_1}
 \otimes v\otimes e^{b_0}\\
 &\hspace{45mm}{}
 \otimes z_1\otimes e^{b_0}\cdots
 e^{b_0}\otimes z_{k_0}\otimes e^{b_0}
 \bigr),
\end{aligned}
\]
其中
\[
 \mathbf x=x_1\otimes\cdots\otimes x_{k_1},
 \qquad
 \mathbf z=z_1\otimes\cdots\otimes z_{k_0}.
\]

\begin{lemma}{3.8.74}\label{lem:3-8-74}
\begin{enumerate}[label=\textup{(\arabic*)},leftmargin=2.8em]
\item 算子 $\mathfrak n^{\mathfrak b}_{k_1,k_0}$ 在
$C_{1+}(L^{(1)},L^{(0)};\Lambda_{0,\mathrm{nov}})$ 上定义一个左、右
代数分别为
\[
 \bigl(C(L^{(1)};\Lambda_{0,\mathrm{nov}}),
       \mathfrak m^{\mathfrak b,0}\bigr),
 \qquad
 \bigl(C(L^{(0)};\Lambda_{0,\mathrm{nov}}),
       \mathfrak m^{\mathfrak b,0}\bigr)
\]
的滤过 $A_\infty$ 双模结构。

\item 算子
$ {}^{(\mathfrak b,b_1)}
\mathfrak n^{(\mathfrak b,b_0)}$
也在 $C_{1+}(L^{(1)},L^{(0)};\Lambda_{0,\mathrm{nov}})$ 上定义一个
左、右代数分别为
\[
 \bigl(C(L^{(1)};\Lambda_{0,\mathrm{nov}}),
       \mathfrak m^{\mathfrak b,b_1}\bigr),
 \qquad
 \bigl(C(L^{(0)};\Lambda_{0,\mathrm{nov}}),
       \mathfrak m^{\mathfrak b,b_0}\bigr)
\]
的滤过 $A_\infty$ 双模结构。
\end{enumerate}
\end{lemma}

\begin{proof}
这是基于 \eqref{eq:3-8-72} 的直接计算。
\end{proof}

引理 3.8.74 表明，$\mathfrak n^{\mathfrak b}_{k_1,k_0}$ 定义了由
$\mathfrak b$ 形变的滤过 $A_\infty$ 双模结构，而
$ {}^{(\mathfrak b,b_1)}
\mathfrak n^{(\mathfrak b,b_0)}$
可视为 bulk 形变后的形变滤过 $A_\infty$ 双模结构。特别地，考虑
\[
\begin{aligned}
 \delta_{(\mathfrak b,b_1),(\mathfrak b,b_0)}
 \colon
 C_{1+}^k(L^{(1)},L^{(0)};\Lambda_{0,\mathrm{nov}})
 \longrightarrow
 C_{1+}^{k+1}(L^{(1)},L^{(0)};\Lambda_{0,\mathrm{nov}})
\end{aligned}
\]
并定义
\[
 \delta_{(\mathfrak b,b_1),(\mathfrak b,b_0)}(v)
 =
 \mathfrak r(
 e^{\mathfrak b}\otimes e^{b_1}\otimes v\otimes e^{b_0}
 ).
\]
于是有：

\begin{statement}{引理--定义}{3.8.75}\label{lem-def:3-8-75}
\begin{enumerate}[label=\textup{(\arabic*)},leftmargin=2.8em]
\item 若
\[
 b_0\in\widehat{\mathcal M}^{\mathfrak b}(L^{(0)}),
 \qquad
 b_1\in\widehat{\mathcal M}^{\mathfrak b}(L^{(1)}),
\]
则
\[
 \delta_{(\mathfrak b,b_1),(\mathfrak b,b_0)}
 \circ
 \delta_{(\mathfrak b,b_1),(\mathfrak b,b_0)}
 =0.
\]
因此可以定义弗洛尔上同调群
\[
 HF\bigl(
 (L^{(1)},(\mathfrak b,b_1)),
 (L^{(0)},(\mathfrak b,b_0));
 \Lambda_{0,\mathrm{nov}}
 \bigr)
\]
为上链复形
\[
 \bigl(
 C_{1+}(L^{(1)},L^{(0)};\Lambda_{0,\mathrm{nov}}),
 \delta_{(\mathfrak b,b_1),(\mathfrak b,b_0)}
 \bigr)
\]
的上同调。

\item 若 $L^{(i)}$（$i=0,1$）经 bulk 形变后弱无障碍（见定义
3.8.40\translatornote{原书此处写“定义 3.8.39”，但 3.8.39 的题头
是引理；经 bulk 形变后弱无障碍及其模空间实际定义于紧随其后的
定义 3.8.40。译文改正题头编号。}），则对任意满足
\[
 (\mathfrak b,b_i)
 \in\widehat{\mathcal M}_{\mathrm{weak,def}}(L^{(i)})
\]
以及
\[
 \mathfrak{PO}_{L^{(0)}}(\mathfrak b,b_0)
 =
 \mathfrak{PO}_{L^{(1)}}(\mathfrak b,b_1)
\]
的 $(\mathfrak b,b_i)$，定义
\[
 \delta^+_{(\mathfrak b,b_1),(\mathfrak b,b_0)}(v)
 =
 \mathfrak r^+\bigl(
 e^{\mathfrak b}\otimes e^{b_1}\otimes v\otimes e^{b_0}
 \bigr).
\]
此时
\[
 \delta^+_{(\mathfrak b,b_1),(\mathfrak b,b_0)}
 \circ
 \delta^+_{(\mathfrak b,b_1),(\mathfrak b,b_0)}
 =0.
\]
因此同样可以定义弗洛尔上同调群
\[
 HF\bigl(
 (L^{(1)},(\mathfrak b,b_1)),
 (L^{(0)},(\mathfrak b,b_0));
 \Lambda_{0,\mathrm{nov}}
 \bigr)
\]
为上链复形
\[
 \bigl(
 C_{1+}(L^{(1)},L^{(0)};\Lambda_{0,\mathrm{nov}}),
 \delta^+_{(\mathfrak b,b_1),(\mathfrak b,b_0)}
 \bigr)
\]
的上同调。
\end{enumerate}
\end{statement}
\liftIndex{I015}
\liftIndex{I125}\liftIndex{I329}

在引理--定义 3.8.75 (2) 中\liftIndex{I200}，还可以类似地定义 bulk 形变后的形变滤过
$A_\infty$ 双模结构：
\[
\begin{aligned}
&{}^{(\mathfrak b,b_1)}
\mathfrak n^{+(\mathfrak b,b_0)}
(\mathbf x\otimes v\otimes\mathbf z)\\
&\quad=
\mathfrak r^+\bigl(
e^{\mathfrak b}\otimes
e^{b_1}\otimes x_1\otimes e^{b_1}\cdots
e^{b_1}\otimes x_{k_1}\otimes e^{b_1}
\otimes v\otimes e^{b_0}\\
&\hspace{43mm}{}
\otimes z_1\otimes e^{b_0}\cdots
e^{b_0}\otimes z_{k_0}\otimes e^{b_0}
\bigr).
\end{aligned}
\]

\subsection{算子 \texorpdfstring{$\mathfrak r$}{r} 的构造}
\label{subsec:3-8-8}
\liftIndex{I201}

本小节讨论定理 3.8.71 的证明。其证明是前述构造的直接组合。取殆复
结构的单参数族 $\{J_t\}_t$。考虑满足下列条件的
$((\Sigma,\vec\tau,\vec z^+),u)$：

\noindent\textup{(3.8.76.1)}
$((\Sigma,\vec\tau),u)$ 满足与模空间
\[
 \mathcal M_{k_1,k_0}
 \bigl(L^{(1)},L^{(0)};[h,w],[h',w']\bigr)
\]
中元素相同的条件，但可以含有不稳定分支。

\noindent\textup{(3.8.76.2)}
$\vec z^+=(z_1^+,\ldots,z_\ell^+)$ 是 $\Sigma$ 上的内部点。

\noindent\textup{(3.8.76.3)}
$((\Sigma,\vec\tau,\vec z^+),u)$ 是稳定的，即其自同构群是有限群。

这里
\[
 \vec\tau=(\vec\tau^{\,(0)},\vec\tau^{\,(1)}),
 \qquad
 \vec\tau^{\,(i)}
 =\bigl(\tau_1^{(i)},\ldots,\tau_{k_i}^{(i)}\bigr)
 \quad(i=0,1),
\]
如第 3.7.5 小节。用
\[
 \mathcal M_{k_1,k_0;\ell}
 \bigl(L^{(1)},L^{(0)};[h,w],[h',w']\bigr)
\]
表示这类组
\[
 ((\Sigma,\vec\tau,\vec z^+),u)
\]
的模空间。定义评价映射
\[
\begin{aligned}
 &(ev_{\mathrm{int}},ev^{(1)},ev^{(0)})\colon
 \mathcal M_{k_1,k_0;\ell}
 \bigl(L^{(1)},L^{(0)};[h,w],[h',w']\bigr)\\
 &\hspace{43mm}\longrightarrow
 M^\ell\times(L^{(1)})^{k_1}\times(L^{(0)})^{k_0}
\end{aligned}
\]
如下：
\[
\begin{aligned}
 &(ev_{\mathrm{int}},ev^{(1)},ev^{(0)})
 \bigl((\Sigma,\vec\tau,\vec z^+),u\bigr)\\
 &\quad=
 \bigl(
 u(z_1^+),\ldots,u(z_\ell^+),
 u((\tau_1^{(1)},1)),\ldots,u((\tau_{k_1}^{(1)},1)),\\
 &\hspace{34mm}
 u((\tau_1^{(0)},0)),\ldots,u((\tau_{k_0}^{(0)},0))
 \bigr).
\end{aligned}
\]
又定义
\[
\begin{aligned}
 ev_{-\infty}\colon&
 \mathcal M_{k_1,k_0;\ell}
 \bigl(L^{(1)},L^{(0)};[h,w],[h',w']\bigr)
 \longrightarrow R_h,\\
 ev_{+\infty}\colon&
 \mathcal M_{k_1,k_0;\ell}
 \bigl(L^{(1)},L^{(0)};[h,w],[h',w']\bigr)
 \longrightarrow R_{h'},
\end{aligned}
\]
其中
\[
 ev_{\pm\infty}
 \bigl((\Sigma,\vec\tau,\vec z^+),u\bigr)
 =\lim_{\tau\to\pm\infty}u(\tau,t).
\]

设
\[
 Q_i\in\mathcal X(M),\qquad
 P_i^{(j)}\in\mathcal X_{1+}(L^{(j)}),\qquad
 S\in\mathcal X_{1+}(R_h).
\]
记
\[
 \mathcal E
 =\bigl(ev_{\mathrm{int}},ev^{(1)},ev_{-\infty},ev^{(0)}\bigr),
\]
并定义
\[
\begin{aligned}
 &\mathcal M_{k_1,k_0;\ell}
 \bigl(L^{(1)},L^{(0)};[h,w],[h',w'];
 \vec Q,\vec P^{(1)},S,\vec P^{(0)}\bigr)\\
 &\quad=
 \mathcal M_{k_1,k_0;\ell}
 \bigl(L^{(1)},L^{(0)};[h,w],[h',w']\bigr)\!
 _{\mathcal E}\!
 \mathbin{\times}
 \left(
 \prod_i Q_i\times\prod_i P_i^{(1)}
 \times S\times\prod_i P_i^{(0)}
 \right).
\end{aligned}
\]
见图 3.8.5。

\begin{figure}[htbp]
  \centering
  \setlength{\unitlength}{0.9mm}
  \begin{picture}(140,64)(0,-5)
    \qbezier(9,30)(62,63)(113,30)
    \qbezier(9,30)(62,-3)(113,30)
    \qbezier[50](13,38)(62,73)(109,38)
    \qbezier[50](13,22)(62,-13)(109,22)
    \put(5,25){\line(1,1){9}}
    \put(5,35){\line(1,-1){9}}
    \put(108,25){\line(1,1){9}}
    \put(108,35){\line(1,-1){9}}
    \put(28,48){\line(1,-2){3}}
    \put(91,48){\line(-1,-2){3}}
    \put(28,12){\line(1,2){3}}
    \put(91,12){\line(-1,2){3}}
    \put(18,50){$P_{k_1}^{(1)}$}
    \put(87,50){$P_1^{(1)}$}
    \put(18,3){$P_1^{(0)}$}
    \put(86,3){$P_{k_0}^{(0)}$}
    \put(1,28){$S$}
    \put(116,30){\vector(1,0){5}}
    \put(122,27){$ev_{+\infty}$}
    \put(46,31){\circle*{1.6}}
    \put(68,38){\circle*{1.6}}
    \put(62,22){\circle*{1.6}}
    \put(46,32){\line(0,1){8}}
    \put(68,39){\line(1,2){4}}
    \put(62,23){\line(0,1){8}}
    \put(44,42){$Q_1$}
    \put(72,46){$Q_2$}
    \put(60,33){$Q_3$}
  \end{picture}
  \par\smallskip
  \renewcommand{\thefigure}{3.8.5}
  \refstepcounter{figure}
  \textbf{\figurename~\thefigure}\label{fig:3-8-5}
\end{figure}

该模空间带有仓西结构。取其上的多值截面 $\mathfrak s$，使之在端部
相容，并置
\[
\begin{aligned}
 &\mathcal M_{k_1,k_0;\ell}
 \bigl(L^{(1)},L^{(0)};[h,w],[h',w'];
 \vec Q,\vec P^{(1)},S,\vec P^{(0)}\bigr)^{\mathfrak s}\\
 &\qquad=
 (-1)^\epsilon ev_{+\infty}
 \bigl(\mathfrak s^{-1}(0)\bigr),
\end{aligned}
\]
这是 $R_{h'}$ 上的一条链。这里的符号 $(-1)^\epsilon$ 与第 3.7.5
小节定义 3.7.71 之前的符号相同。可以选择
$\mathcal X_{1+}(R_h)$，使虚基本链
\[
 \mathcal M_{k_1,k_0;\ell}
 \bigl(L^{(1)},L^{(0)};[h,w],[h',w'];
 \vec Q,\vec P^{(1)},S,\vec P^{(0)}\bigr)^{\mathfrak s}
\]
仍包含于其中。现在取下列算子\translatornote{原书此处把算子写成
$\mathfrak r_{k_1,k_0;\ell}$，但定理 3.8.71 从定义域起始即固定记号
$\mathfrak r_{\ell;k_1,k_0}$，其各恒等式也全都沿用这一指标次序；
定理 7.4.154 的证明再次明确调用
$\mathfrak r_{\ell;k_1,k_0}$。译文依该全局记号恢复指标次序。}：
\[
\begin{aligned}
 &\mathfrak r_{\ell;k_1,k_0}
 \bigl(
 Q_1\otimes\cdots\otimes Q_\ell
 \otimes P_1^{(1)}\otimes\cdots\otimes P_{k_1}^{(1)}
 \otimes[S,w]\\
 &\hspace{35mm}{}
 \otimes P_1^{(0)}\otimes\cdots\otimes P_{k_0}^{(0)}
 \bigr)\\
 &\quad=
 \sum_{[h',w']}\frac1{\ell!}
 \mathcal M_{k_1,k_0;\ell}
 \bigl(L^{(1)},L^{(0)};[h,w],[h',w'];
 \vec Q,\vec P^{(1)},S,\vec P^{(0)}\bigr)^{\mathfrak s}.
\end{aligned}
\]
式 \eqref{eq:3-8-72} 与 (3.8.73) 的证明由此已显然。
